\documentclass[UTF8,12pt,a4paper]{ctexart}

\usepackage{amsmath}
\usepackage{amssymb}
\usepackage{graphicx}
\usepackage{enumerate}
\usepackage[margin=1in]{geometry}
\geometry{a4paper}
\usepackage{fancyhdr}
\pagestyle{fancy}
\fancyhf{}
\usepackage{xcolor}
\usepackage{hyperref}
\hypersetup{
    pdfborder={0 0 0},
    colorlinks=true,
    linkcolor=black,
    citecolor=black,
    urlcolor=blue
}
\usepackage{booktabs}
\usepackage{array}
\usepackage{multirow}
\usepackage{longtable}
\newcolumntype{L}[1]{>{\raggedright\arraybackslash}p{#1}}
\newcolumntype{M}[1]{>{\centering\arraybackslash}p{#1}}

\title{集成电路工艺技术基础 Chapter 2\\晶圆制造与清洗详细复习手稿}
\author{冯峻}
\date{\today}

\fancyhead[L]{Chapter 2 复习手稿}
\fancyhead[C]{晶圆制造与清洗}
\fancyfoot[C]{\thepage}

\begin{document}

\maketitle
\tableofcontents
\newpage

\section{硅晶体结构与缺陷}

\subsection{硅晶体基本结构}

\subsubsection{晶体结构特征}
\begin{itemize}
    \item \textbf{晶格类型：}金刚石结构（Diamond structure）
    \item \textbf{晶格常数：}a = 5.43 \AA（埃）
    \item \textbf{原子密度：}5 $\times$ 10$^{22}$ atoms/cm$^3$
    \item \textbf{键合方式：}共价键，每个Si原子与4个相邻Si原子形成sp$^3$杂化轨道
    \item \textbf{晶体取向：}常用(100)、(110)、(111)晶面，CMOS工艺主要使用(100)晶面
\end{itemize}

\subsubsection{为什么选择硅？}
\begin{enumerate}
    \item 地壳中储量丰富（第二丰富元素，仅次于氧）
    \item 可形成高质量SiO$_2$介电层（Si/SiO$_2$界面质量优异）
    \item 室温下带隙适中（E$_g$ = 1.12 eV），适合电子器件
    \item 机械强度好，易于加工
    \item 工艺技术成熟
\end{enumerate}

\subsection{晶体缺陷分类}

晶体缺陷对器件性能有重要影响，必须严格控制。

\subsubsection{点缺陷（Zero-dimensional defects）}

点缺陷主导大多数掺杂扩散机制，决定杂质分布轮廓。

\begin{enumerate}
    \item \textbf{替位式原子（Substitutional）}
    \begin{itemize}
        \item 杂质原子占据正常晶格位置
        \item 例：掺杂原子B、P、As等
        \item 是主要的电学激活掺杂形式
    \end{itemize}
    
    \item \textbf{间隙式原子（Interstitial）}
    \begin{itemize}
        \item 原子位于晶格间隙位置
        \item 扩散速度快
        \item 某些掺杂原子以间隙式存在时不具有电学活性
    \end{itemize}
    
    \item \textbf{空位（Vacancy）}
    \begin{itemize}
        \item 晶格位置缺少原子
        \item 参与扩散过程
    \end{itemize}
    
    \item \textbf{肖特基缺陷（Schottky defect）}
    \begin{itemize}
        \item 表面原子空位
    \end{itemize}
    
    \item \textbf{弗伦克尔对（Frenkel-type interstitial-vacancy pair）}
    \begin{itemize}
        \item 间隙原子-空位对
    \end{itemize}
\end{enumerate}

\subsubsection{线缺陷（One-dimensional defects）}

\textbf{位错（Dislocation）：}缺失或额外的原子线

\begin{enumerate}
    \item \textbf{刃位错（Edge dislocation）}
    \begin{itemize}
        \item 由多余或缺失的半原子面形成
        \item 垂直于位错线方向
    \end{itemize}
    
    \item \textbf{螺位错（Screw dislocation）}
    \begin{itemize}
        \item 使晶面呈现螺旋结构
        \item 常在晶体生长中出现
        \item 平行于位错线方向
    \end{itemize}
\end{enumerate}

\textbf{位错的危害：}
\begin{itemize}
    \item 降低载流子迁移率
    \item 增加漏电流
    \item 降低击穿电压
    \item 影响器件可靠性
\end{itemize}

\subsubsection{面缺陷（Two-dimensional defects）}

\begin{enumerate}
    \item \textbf{层错（Stacking fault）}
    \begin{itemize}
        \item 晶体层堆垛顺序的破坏
        \item 由位错终止
        \item 影响载流子输运
    \end{itemize}
    
    \item \textbf{晶界（Grain boundary）}
    \begin{itemize}
        \item 不同晶粒之间的界面
        \item 单晶硅中应避免
    \end{itemize}
    
    \item \textbf{表面（Surface）}
    \begin{itemize}
        \item 晶体与外界的界面
        \item 表面态对器件性能影响显著
    \end{itemize}
\end{enumerate}

\subsubsection{体缺陷（Three-dimensional defects）}

\textbf{析出物（Precipitation）}

\begin{itemize}
    \item \textbf{形成原因：}杂质浓度超过固溶度限，冷却后形成过饱和状态，过量杂质析出成第二相
    \item \textbf{常见类型：}硅中的氧析出物（SiO$_2$）
    \item \textbf{作用：}
    \begin{itemize}
        \item 正面：内吸杂作用，捕获有害金属杂质
        \item 负面：过大析出物可能成为缺陷源
    \end{itemize}
\end{itemize}

\section{单晶硅生长}

\subsection{原材料制备}

\subsubsection{从沙子到冶金级硅（MGS）}

\textbf{化学反应：}
\begin{equation}
2\text{SiO}_2 + 2\text{C} \rightarrow \text{Si(liquid)} + 2\text{CO}
\end{equation}

\begin{itemize}
    \item \textbf{原料：}天然砂（主要成分SiO$_2$）
    \item \textbf{产物：}冶金级硅（Metallurgical Grade Silicon, MGS）
    \item \textbf{纯度：}约98-99\%，含有较多杂质
    \item \textbf{用途：}需要进一步提纯才能用于半导体制造
\end{itemize}

\subsubsection{从MGS到电子级硅（EGS）}

\textbf{三氯氢硅法（Siemens工艺）：}

\textbf{化学反应：}
\begin{equation}
\text{Si} + 3\text{HCl} \leftrightarrow \text{SiHCl}_3\text{(liquid)} + \text{H}_2
\end{equation}

\textbf{提纯过程：}
\begin{enumerate}
    \item MGS与HCl反应生成三氯氢硅（SiHCl$_3$）液体
    \item 通过蒸馏纯化SiHCl$_3$（去除杂质）
    \item 将纯化的SiHCl$_3$在高温下还原沉积成高纯硅
\end{enumerate}

\textbf{电子级硅（EGS）纯度指标：}
\begin{itemize}
    \item 总杂质：< 1 ppb（约5 $\times$ 10$^{13}$ cm$^{-3}$）
    \item 氧含量：约10$^{18}$ cm$^{-3}$（来自工艺过程）
    \item 碳含量：约10$^{16}$ cm$^{-3}$（来自工艺过程）
    \item \textbf{质量标准：}达到半导体制造要求
\end{itemize}

\subsection{Czochralski（CZ）法单晶生长}

\subsubsection{历史背景}

\begin{itemize}
    \item \textbf{发明者：}Jan Czochralski
    \item \textbf{发明时间：}1916年
    \item \textbf{发现过程：}Czochralski意外将笔浸入熔融锡坩埚（原本想蘸墨水），立即拉出发现笔尖悬挂一根固化金属丝，验证后发现是单晶
    \item \textbf{应用：}目前几乎所有IC用硅片都采用CZ法生长
\end{itemize}

\subsubsection{CZ法工艺原理}

\textbf{基本流程：}

\begin{enumerate}
    \item \textbf{装料：}将EGS硅破碎成小块，放入SiO$_2$（石英）坩埚
    \item \textbf{加热熔化：}在Ar气氛保护下，通过RF（射频）线圈加热坩埚
    \item \textbf{温度控制：}加热至略高于1417°C（硅熔点）
    \item \textbf{掺杂：}向熔体中加入掺杂剂（B、P、As等），控制导电类型和电阻率
    \item \textbf{引晶：}单晶籽晶（seed crystal）浸入熔体
    \item \textbf{提拉生长：}
    \begin{itemize}
        \item 籽晶缓慢向上提拉
        \item 熔体沿籽晶上升并冷却结晶
        \item 同时旋转籽晶（抑制各向异性生长，获得圆柱形晶体）
    \end{itemize}
    \item \textbf{形成晶锭：}生长出长圆柱形单晶锭（boule/ingot）
\end{enumerate}

\textbf{关键控制参数：}
\begin{itemize}
    \item \textbf{提拉速度：}控制生长速率和晶体质量
    \item \textbf{熔体温度：}影响生长界面形状和缺陷形成
    \item \textbf{旋转速度：}保证圆形截面和均匀生长
    \item \textbf{温度梯度：}影响位错密度
\end{itemize}

\subsubsection{CZ法设备构成}

\begin{itemize}
    \item \textbf{籽晶夹持器（Seed holder）：}固定并旋转籽晶
    \item \textbf{石英坩埚（Quartz crucible）：}盛装熔融硅
    \item \textbf{石墨托架（Graphite susceptor）：}支撑坩埚
    \item \textbf{RF加热线圈：}感应加热
    \item \textbf{Ar气氛保护：}防止氧化
    \item \textbf{固-液界面：}晶体生长的关键区域
\end{itemize}

\subsubsection{CZ法产品特性}

\begin{itemize}
    \item \textbf{晶锭重量：}约100 kg
    \item \textbf{晶锭长度：}可达1-2米
    \item \textbf{直径：}目前主流300 mm（12英寸），最新450 mm
    \item \textbf{晶向：}通常为<100>方向
    \item \textbf{圆形截面：}通过旋转控制获得
\end{itemize}

\subsubsection{CZ硅中的氧和碳}

\textbf{来源：}
\begin{itemize}
    \item \textbf{氧（O）：}来自SiO$_2$坩埚，高温下部分溶解到硅熔体中
    \item \textbf{碳（C）：}来自石墨托架/加热器
\end{itemize}

\textbf{典型浓度：}
\begin{itemize}
    \item C$_O$ $\approx$ 10$^{18}$ cm$^{-3}$
    \item C$_C$ $\approx$ 10$^{16}$ cm$^{-3}$
\end{itemize}

\textbf{氧的作用：}
\begin{enumerate}
    \item \textbf{机械强度：}提高硅片机械强度，减少翘曲
    \item \textbf{内吸杂（Internal Gettering）：}
    \begin{itemize}
        \item 在正常工艺温度循环下，O形成SiO$_2$析出物
        \item 析出物及其周围缺陷可捕获有害金属杂质（Au、Cu、Fe等）
        \item 将杂质"吸"到远离器件有源区的体内
        \item 提高器件性能和良率
    \end{itemize}
\end{enumerate}

\subsection{浮区（Float Zone, FZ）法}

\subsubsection{FZ法原理}

\begin{itemize}
    \item \textbf{方法：}通过局部加热形成熔区（molten zone）
    \item \textbf{过程：}熔区沿硅棒移动，实现区域提纯和单晶生长
    \item \textbf{特点：}无坩埚接触，避免了坩埚引入的杂质
\end{itemize}

\subsubsection{FZ法优缺点}

\textbf{优点：}
\begin{itemize}
    \item \textbf{高纯度：}无坩埚污染，O和C含量极低
    \item \textbf{高电阻率：}适合制备高电阻率材料
    \item \textbf{少缺陷：}位错密度低
\end{itemize}

\textbf{缺点：}
\begin{itemize}
    \item 成本较高
    \item 直径受限（一般≤200 mm）
    \item 机械强度较低（缺少O的强化作用）
\end{itemize}

\subsubsection{应用领域}

\begin{itemize}
    \item 功率器件（Power devices）
    \item 聚光太阳能电池（Concentrator solar cells）
    \item 对纯度要求极高的特殊应用
\end{itemize}

\section{晶圆制造工艺}

\subsection{直径研磨（Grinding）}

\textbf{必要性：}
\begin{itemize}
    \item 单晶生长过程中直径存在偏差
    \item 需要获得精确的目标直径
\end{itemize}

\textbf{工艺：}
\begin{itemize}
    \item \textbf{方法：}机械研磨（滚磨/barrelling）
    \item \textbf{目标：}将晶锭研磨至目标直径（如300 mm）
    \item \textbf{要求：}直径均匀，偏差控制在数十微米内
\end{itemize}

\subsection{晶向确认与主平面形成}

\subsubsection{晶向确认}

\textbf{方法：}X射线衍射（X-ray diffraction）

\textbf{目的：}
\begin{itemize}
    \item 确认目标晶向（如<100>方向）
    \item 定位晶面方向
\end{itemize}

\subsubsection{主平面/缺口形成}

\textbf{Flat或Notch标识：}

在晶锭上切割参考平面或缺口，标识：
\begin{enumerate}
    \item \textbf{晶向信息：}主flat平行于某个晶向
    \item \textbf{掺杂类型：}通过flat的数量和长度区分N型/P型
\end{enumerate}

\textbf{常见模式：}
\begin{itemize}
    \item P型(100)：一个大flat + 一个小flat
    \item N型(100)：一个大flat
    \item 现代大尺寸晶圆：使用notch（V形缺口）代替flat，减少材料损失
\end{itemize}

\subsection{切片（Slicing/Sawing）}

\textbf{工艺：}
\begin{itemize}
    \item \textbf{设备：}金刚石切割器（Diamond saw）或线锯（Wire saw）
    \item \textbf{方法：}将晶锭切成薄片
\end{itemize}

\textbf{参数：}
\begin{itemize}
    \item \textbf{厚度：}约700-800 $\mu$m（取决于直径）
    \item \textbf{切割损耗：}尽量减少kerf loss（切缝损耗）
\end{itemize}

\textbf{要求：}
\begin{itemize}
    \item 切割面平整
    \item 表面损伤层最小
    \item 厚度均匀
\end{itemize}

\subsection{晶圆抛光（Polishing）}

\subsubsection{抛光必要性}

\begin{itemize}
    \item 切片后表面粗糙，有损伤层
    \item IC制造要求无损伤的镜面级光滑表面
    \item 表面粗糙度需控制在原子级别
\end{itemize}

\subsubsection{化学机械抛光（CMP）}

\textbf{Chemical-Mechanical Polishing原理：}
\begin{itemize}
    \item 化学作用：抛光液中的化学物质与硅反应，软化表面
    \item 机械作用：研磨颗粒去除软化层
    \item 两者协同作用，实现高质量抛光
\end{itemize}

\textbf{控制参数：}
\begin{itemize}
    \item \textbf{曲率（Bow/Warp）：}晶圆平整度
    \item \textbf{粗糙度（Roughness）：}表面光滑程度
    \item \textbf{总厚度变化（TTV）：}厚度均匀性
\end{itemize}

\subsection{激光标记（Marking）}

\textbf{方法：}激光刻蚀

\textbf{内容：}
\begin{itemize}
    \item 晶圆编号
    \item 生产厂商信息
    \item 批次号
    \item 晶向和掺杂类型
\end{itemize}

\textbf{位置：}晶圆边缘区域

\subsection{晶圆评估与检测}

\subsubsection{晶向检测}

\textbf{方法：}通过flat或notch图案识别

\textbf{信息：}
\begin{itemize}
    \item 掺杂类型（N型/P型）
    \item 晶面方向（(100)、(111)等）
    \item 晶向定位
\end{itemize}

\subsubsection{电阻率测试}

\textbf{四探针法（Four-point probe）：}

\textbf{原理：}
\begin{itemize}
    \item 外侧两个探针通电流I
    \item 内侧两个探针测电压V
    \item 计算电阻率$\rho$ = k $\cdot$ (V/I)，k为几何因子
\end{itemize}

\textbf{目的：}
\begin{itemize}
    \item 确认掺杂浓度
    \item 检查均匀性
    \item 验证是否满足设计要求
\end{itemize}

\subsubsection{缺陷检测}

\textbf{检测项目：}
\begin{itemize}
    \item 点缺陷密度
    \item 位错密度
    \item 层错
    \item 金属污染
    \item 颗粒污染
\end{itemize}

\textbf{检测方法：}
\begin{itemize}
    \item 光学显微镜
    \item 扫描电镜（SEM）
    \item X射线形貌术
    \item 表面粒子扫描
\end{itemize}

\section{晶圆清洗}

\subsection{清洗的必要性}

半导体器件对污染极其敏感，微量污染即可导致器件失效。

\subsubsection{主要污染物类型及危害}

\begin{enumerate}
    \item \textbf{有机物/光刻胶}
    \begin{itemize}
        \item \textbf{危害：}
        \begin{itemize}
            \item 增加栅漏电流
            \item 阻碍后续清洗（掩盖下层污染）
            \item 影响工艺粘附性
        \end{itemize}
    \end{itemize}
    
    \item \textbf{颗粒污染}
    \begin{itemize}
        \item \textbf{危害：}
        \begin{itemize}
            \item 致命缺陷（killer defects）：导致器件短路
            \item 栅氧化层可靠性问题
            \item 可能携带金属杂质
        \end{itemize}
        \item \textbf{严重性：}随工艺节点缩小，颗粒尺寸要求越来越严格
    \end{itemize}
    
    \item \textbf{碱金属离子（Alkali metals）}
    \begin{itemize}
        \item \textbf{主要元素：}Na$^+$、K$^+$、Ca$^{2+}$等
        \item \textbf{危害：}
        \begin{itemize}
            \item 在氧化层中移动性强
            \item 响应外加电场
            \item 引起阈值电压漂移
            \item 降低氧化层可靠性
        \end{itemize}
        \item \textbf{定量影响：}
        \begin{itemize}
            \item 栅氧厚度t$_{ox}$ = 10 nm
            \item 0.1V的V$_{th}$漂移对应Q$_M$ = 2.2$\times$10$^{11}$ cm$^{-2}$
            \item 这仅相当于氧化层中<0.1\%的单层污染
        \end{itemize}
    \end{itemize}
    
    \item \textbf{重金属/过渡金属}
    \begin{itemize}
        \item \textbf{主要元素：}Au、Cu、Fe、Ni等
        \item \textbf{危害：}
        \begin{itemize}
            \item 扩散速度快，在温度步骤中快速进入晶圆
            \item 形成深能级陷阱
            \item 成为复合中心，降低载流子寿命
            \item 降低载流子迁移率
            \item 增加漏电流
        \end{itemize}
        \item \textbf{定量要求（以DRAM为例）：}
        \begin{itemize}
            \item 需要刷新时间 > 25 $\mu$s
            \item 要求生成寿命$\tau_G$ > 25 $\mu$s
            \item 这对金属杂质浓度提出极严格要求
        \end{itemize}
        \item \textbf{特别注意：}Au、Cu、Fe等倾向于在晶圆表面和结区富集，危害更大
    \end{itemize}
\end{enumerate}

\subsubsection{污染控制策略}

\begin{itemize}
    \item 洁净室环境（Clean room facilities）
    \item 清洗工艺（Cleaning processes）
    \item 两者必须结合，缺一不可
\end{itemize}

\subsection{洁净室环境}

\subsubsection{洁净室等级（Clean room Class）}

\textbf{定义：}每立方英尺（cubic foot, 1 ft$^3$ = 0.03 m$^3$）空气中粒径$\geq$ 0.5 $\mu$m的颗粒数

\textbf{常见等级：}
\begin{itemize}
    \item \textbf{Class 1：}$\leq$ 1个颗粒/ft$^3$（最高等级）
    \item \textbf{Class 10：}$\leq$ 10个颗粒/ft$^3$
    \item \textbf{Class 100：}$\leq$ 100个颗粒/ft$^3$
    \item \textbf{Class 1000：}$\leq$ 1000个颗粒/ft$^3$
\end{itemize}

\textbf{工业fab：}通常为Class 1-10

\textbf{大学实验室：}通常为Class 100-1000

\subsubsection{洁净室设施}

\begin{enumerate}
    \item \textbf{HEPA过滤器}
    \begin{itemize}
        \item High Efficiency Particulate Air filter
        \item 高效颗粒空气过滤器
        \item 可去除99.97\%以上的0.3 $\mu$m颗粒
    \end{itemize}
    
    \item \textbf{空气循环系统}
    \begin{itemize}
        \item 层流（Laminar flow）设计
        \item 空气从天花板向下流动
        \item 从地板排出并循环过滤
    \end{itemize}
    
    \item \textbf{人员防护}
    \begin{itemize}
        \item 防尘服（"Bunny suits"或"Space suits"）
        \item 覆盖全身，包括头部、手部
        \item 人是最大的颗粒源
    \end{itemize}
    
    \item \textbf{化学品、水和气体过滤}
    \begin{itemize}
        \item 去离子水（DI water）：电阻率>18 M$\Omega\cdot$cm
        \item 高纯化学品：ppb级杂质
        \item 高纯气体：过滤后使用
    \end{itemize}
\end{enumerate}

\subsection{清洗工艺}

\subsubsection{清洗分类}

\begin{enumerate}
    \item \textbf{前道清洗（Front-end cleaning）}
    \begin{itemize}
        \item \textbf{时机：}每次高温工艺步骤之前
        \item \textbf{化学品：}可使用强腐蚀性溶液（此时晶圆上无金属层）
        \item \textbf{目的：}彻底去除各类污染
    \end{itemize}
    
    \item \textbf{后道清洗（Back-end cleaning）}
    \begin{itemize}
        \item \textbf{时机：}金属布线之后
        \item \textbf{化学品：}使用温和溶液
        \item \textbf{限制：}避免腐蚀金属层（Al、Cu等）
    \end{itemize}
\end{enumerate}

\subsubsection{有机物清洗}

\textbf{目的：}去除光刻胶、有机污染物

\textbf{方法1：干法清洗}
\begin{itemize}
    \item \textbf{工艺：}氧等离子体（O$_2$ plasma）
    \item \textbf{反应：}O$_2$ + 有机物 $\rightarrow$ CO$_2$ + H$_2$O
    \item \textbf{优点：}环保、无液体废弃物
\end{itemize}

\textbf{方法2：湿法清洗（Piranha清洗）}
\begin{itemize}
    \item \textbf{配方：}H$_2$SO$_4$(98\%) : H$_2$O$_2$(30\%) = 2:1 至 4:1
    \item \textbf{温度：}热溶液（80-120°C）
    \item \textbf{反应：}强氧化作用，分解有机物为CO$_2$和H$_2$O
    \item \textbf{注意：}反应放热剧烈，需小心操作
\end{itemize}

\subsubsection{RCA清洗（标准清洗流程）}

\textbf{开发历史：}
\begin{itemize}
    \item \textbf{开发者：}Werner Kern
    \item \textbf{时间：}1965年
    \item \textbf{机构：}RCA公司（Radio Corporation of America）
    \item \textbf{地位：}半导体工艺的标准清洗流程，沿用至今
\end{itemize}

\textbf{RCA清洗完整流程：}

\begin{enumerate}
    \item \textbf{SC-1清洗（APM: Ammonia Peroxide Mixture）}
    
    \textbf{配方：}NH$_4$OH : H$_2$O$_2$ : H$_2$O = 1:1:5
    
    \textbf{条件：}70-80°C，10分钟
    
    \textbf{目的：}
    \begin{itemize}
        \item 去除有机物
        \item 去除颗粒污染
    \end{itemize}
    
    \textbf{机理：}
    \begin{itemize}
        \item NH$_4$OH提供碱性环境
        \item H$_2$O$_2$氧化有机物
        \item 轻微腐蚀硅表面（约1 nm/min），带走颗粒
        \item 表面带负电，颗粒也带负电，静电排斥有助于去除颗粒
    \end{itemize}
    
    \item \textbf{稀HF浸泡}
    
    \textbf{配方：}HF : H$_2$O = 1:50（稀释）
    
    \textbf{条件：}室温，1分钟
    
    \textbf{目的：}去除SC-1清洗后生成的薄氧化层
    
    \textbf{反应：}
    \begin{equation}
    \text{SiO}_2 + 6\text{HF} \rightarrow \text{H}_2\text{SiF}_6 + 2\text{H}_2\text{O}
    \end{equation}
    
    \textbf{结果：}获得氢钝化的硅表面（Si-H）
    
    \item \textbf{SC-2清洗（HPM: Hydrochloric Peroxide Mixture）}
    
    \textbf{配方：}HCl : H$_2$O$_2$ : H$_2$O = 1:1:6
    
    \textbf{条件：}70-80°C，10分钟
    
    \textbf{目的：}
    \begin{itemize}
        \item 去除碱金属离子（Al$^{3+}$、Fe$^{3+}$、Mg$^{2+}$等）
        \item 去除重金属（Au、Cu等）
    \end{itemize}
    
    \textbf{机理：}
    \begin{itemize}
        \item HCl提供酸性环境和Cl$^-$
        \item 金属离子与Cl$^-$形成可溶性氯化物
        \item H$_2$O$_2$氧化金属，促进溶解
    \end{itemize}
    
    \textbf{典型反应：}
    \begin{equation}
    \text{Fe} + 2\text{HCl} + \text{H}_2\text{O}_2 \rightarrow \text{FeCl}_2 + 2\text{H}_2\text{O}
    \end{equation}
\end{enumerate}

\subsubsection{现代RCA清洗改进}

\begin{itemize}
    \item 降低化学品浓度（减少硅损耗和表面粗糙度）
    \item 优化温度和时间
    \item 添加表面活性剂（增强颗粒去除）
    \item 采用兆声波（Megasonic）辅助清洗
\end{itemize}

\subsection{ITRS对污染的要求趋势}

\textbf{随着技术节点缩小：}
\begin{itemize}
    \item 允许的颗粒尺寸越来越小
    \item 允许的颗粒密度越来越低
    \item 金属污染浓度要求越来越严格
    \item 清洗工艺越来越具有挑战性
\end{itemize}

\section{章节总结与知识点梳理}

\subsection{核心概念对照}

\begin{longtable}{|L{4cm}|L{11cm}|}
\hline
\textbf{概念} & \textbf{关键要点} \\
\hline
\endfirsthead
\hline
\textbf{概念} & \textbf{关键要点} \\
\hline
\endhead
\hline
\endfoot
\hline
\endlastfoot

硅晶体结构 & 金刚石结构，a=5.43\AA，5$\times$10$^{22}$ atoms/cm$^3$，共价键，sp$^3$杂化 \\
\hline

点缺陷 & 替位式、间隙式、空位、肖特基、弗伦克尔对；主导扩散机制 \\
\hline

线缺陷 & 刃位错、螺位错；降低迁移率，增加漏电 \\
\hline

面缺陷 & 层错、晶界、表面；影响载流子输运 \\
\hline

体缺陷 & 析出物；O析出可实现内吸杂 \\
\hline

MGS制备 & 2SiO$_2$ + 2C $\rightarrow$ Si + 2CO；纯度98-99\% \\
\hline

EGS提纯 & 三氯氢硅法；杂质<1ppb；O$\sim$10$^{18}$，C$\sim$10$^{16}$cm$^{-3}$ \\
\hline

CZ法 & Czochralski 1916年发明；籽晶提拉旋转生长；主流工艺 \\
\hline

CZ法O和C & O来自SiO$_2$坩埚，C来自石墨；O提供强度和内吸杂 \\
\hline

FZ法 & 浮区法；无坩埚；高纯度；用于功率器件 \\
\hline

直径研磨 & 机械研磨至目标直径；补偿生长偏差 \\
\hline

晶向确认 & X射线衍射；形成flat/notch标识晶向和掺杂类型 \\
\hline

切片 & 金刚石切割；厚度700-800$\mu$m；最小化损伤 \\
\hline

CMP抛光 & 化学机械协同；获得无损伤镜面；控制曲率和粗糙度 \\
\hline

激光标记 & 晶圆编号、厂商信息、批次号 \\
\hline

电阻率测试 & 四探针法；确认掺杂浓度和均匀性 \\
\hline

污染危害 & 有机物、颗粒、碱金属、重金属；导致V$_{th}$漂移、寿命降低、漏电 \\
\hline

洁净室等级 & Class 1-1000；颗粒数/ft$^3$；HEPA过滤、防尘服 \\
\hline

有机物清洗 & O$_2$等离子体或Piranha(H$_2$SO$_4$:H$_2$O$_2$) \\
\hline

SC-1(APM) & NH$_4$OH:H$_2$O$_2$:H$_2$O=1:1:5；70-80°C；去有机物和颗粒 \\
\hline

稀HF浸泡 & HF:H$_2$O=1:50；室温；去氧化层 \\
\hline

SC-2(HPM) & HCl:H$_2$O$_2$:H$_2$O=1:1:6；70-80°C；去金属离子 \\
\hline

清洗分类 & 前道清洗：强溶液，高温前；后道清洗：温和溶液，金属后 \\
\hline

\end{longtable}

\subsection{重要公式和数据}

\begin{enumerate}
    \item 硅晶格常数：a = 5.43 \AA
    \item 硅原子密度：5 $\times$ 10$^{22}$ cm$^{-3}$
    \item 硅熔点：1417°C
    \item 硅带隙：E$_g$ = 1.12 eV (室温)
    \item EGS纯度：< 1 ppb ($\sim$ 5$\times$10$^{13}$ cm$^{-3}$)
    \item CZ硅O浓度：$\sim$ 10$^{18}$ cm$^{-3}$
    \item CZ硅C浓度：$\sim$ 10$^{16}$ cm$^{-3}$
    \item 晶圆厚度：700-800 $\mu$m
    \item V$_{th}$漂移：t$_{ox}$=10nm时，0.1V漂移对应Q$_M$=2.2$\times$10$^{11}$cm$^{-2}$
    \item SC-1配方：NH$_4$OH:H$_2$O$_2$:H$_2$O = 1:1:5
    \item SC-2配方：HCl:H$_2$O$_2$:H$_2$O = 1:1:6
\end{enumerate}

\subsection{关键工艺流程图}

\textbf{完整晶圆制造流程：}
\begin{enumerate}
    \item 天然砂 $\rightarrow$ MGS（2SiO$_2$ + 2C $\rightarrow$ Si + 2CO）
    \item MGS $\rightarrow$ EGS（三氯氢硅法提纯）
    \item EGS $\rightarrow$ 单晶硅锭（CZ法或FZ法）
    \item 直径研磨
    \item 晶向确认与主平面形成
    \item 切片
    \item CMP抛光
    \item 激光标记
    \item 晶圆评估（晶向、电阻率、缺陷检测）
\end{enumerate}

\textbf{标准清洗流程：}
\begin{enumerate}
    \item 有机物清洗（Piranha或O$_2$等离子体）
    \item SC-1清洗（去有机物和颗粒）
    \item DI水冲洗
    \item 稀HF浸泡（去氧化层）
    \item DI水冲洗
    \item SC-2清洗（去金属离子）
    \item DI水冲洗
    \item 旋转甩干或N$_2$吹干
\end{enumerate}

\subsection{常见考点和易错点}

\begin{enumerate}
    \item \textbf{区分CZ法和FZ法}
    \begin{itemize}
        \item CZ法：主流，有坩埚，O和C含量高，内吸杂
        \item FZ法：高纯，无坩埚，O和C低，用于功率器件
    \end{itemize}
    
    \item \textbf{O在CZ硅中的双重作用}
    \begin{itemize}
        \item 正面：机械强度、内吸杂
        \item 要控制浓度和分布
    \end{itemize}
    
    \item \textbf{RCA清洗各步骤功能}
    \begin{itemize}
        \item SC-1：去有机物和颗粒（碱性）
        \item HF：去氧化层
        \item SC-2：去金属离子（酸性）
        \item 顺序不可颠倒
    \end{itemize}
    
    \item \textbf{污染物的定量影响}
    \begin{itemize}
        \item 必须记住V$_{th}$漂移与Q$_M$的关系
        \item 理解ppb级污染的严重性
    \end{itemize}
    
    \item \textbf{晶向标识}
    \begin{itemize}
        \item Flat模式：大flat+小flat（P型），大flat（N型）
        \item Notch：现代大尺寸晶圆
        \item 目的：工艺对准和掺杂类型识别
    \end{itemize}
\end{enumerate}

\subsection{与后续章节的联系}

\begin{itemize}
    \item \textbf{氧化：}清洗后的硅表面是氧化的起点
    \item \textbf{扩散：}点缺陷机制在扩散中起关键作用
    \item \textbf{离子注入：}注入后需要退火，会影响O析出
    \item \textbf{薄膜沉积：}需要清洁的表面确保粘附性
    \item \textbf{光刻：}光刻胶去除需要清洗工艺
    \item \textbf{刻蚀：}刻蚀后需要清洗去除残留
\end{itemize}

\section{复习建议}

\subsection{重点掌握内容}

\begin{enumerate}
    \item 硅晶体结构和四种缺陷类型（点、线、面、体）
    \item EGS制备的两步反应式
    \item CZ法的原理、流程、关键参数
    \item CZ硅中O和C的来源、浓度、作用
    \item 晶圆制造的完整流程（7个主要步骤）
    \item 四种主要污染物及其危害
    \item RCA清洗的三个步骤（配方、条件、目的、机理）
    \item 洁净室等级定义
\end{enumerate}

\subsection{理解性内容}

\begin{enumerate}
    \item 为什么选择硅作为主要半导体材料？
    \item 内吸杂的原理和意义
    \item CZ法与FZ法的区别和应用场景
    \item 为什么微量污染就能导致器件失效？
    \item SC-1和SC-2为什么使用不同的酸碱性？
    \item 前道清洗和后道清洗的区别
\end{enumerate}

\subsection{计算题准备}

\begin{enumerate}
    \item 给定V$_{th}$漂移量，计算允许的移动离子浓度
    \item 给定载流子寿命要求，推算允许的金属杂质浓度
    \item 洁净室等级的颗粒浓度换算
    \item 电阻率与掺杂浓度的关系
\end{enumerate}

\subsection{记忆技巧}

\begin{itemize}
    \item \textbf{RCA清洗顺序：}碱性去颗粒(SC-1) $\rightarrow$ 去氧化层(HF) $\rightarrow$ 酸性去金属(SC-2)
    \item \textbf{缺陷维度：}点(0D) $\rightarrow$ 线(1D) $\rightarrow$ 面(2D) $\rightarrow$ 体(3D)
    \item \textbf{硅提纯：}沙子 $\rightarrow$ MGS（冶金级）$\rightarrow$ EGS（电子级）
    \item \textbf{CZ vs FZ：}CZ有坩埚、O高、主流；FZ无坩埚、纯净、特殊用途
\end{itemize}

\end{document}
